\subsection{RELATIONS BETWEEN THE VARIOUS FREE OBJECTS}

As the free monoid \(\mathrm{Mo}(\mathbf{X})\) is a magma, Proposition 1, of no. 1 shows the existence of a homomorphism \(\lambda: \mathrm{M}(\mathbf{X}) \to \mathrm{Mo}(\mathbf{X})\) whose restriction to \(\mathbf{X}\) is the identity. Similarly, as the free group \(\mathrm{F}(\mathbf{X})\) is a monoid the identity mapping of \(\mathbf{X}\) extends to a homomorphism \(\mu: \mathrm{Mo}(\mathbf{X}) \to \mathrm{F}(\mathbf{X})\) (no. 2, Proposition 3). By no. 4, Proposition 6 and no. 5, Proposition 7, \(\mu\) is injective. Similarly Proposition 10 of no. 7 and Proposition 8 of no. 5 show the existence of homomorphisms \(\nu: \mathrm{Mo}(\mathbf{X}) \to \mathbf{N}^{(\mathbf{X})}\) and \(\pi: \mathrm{F}(\mathbf{X}) \to \mathbf{Z}^{(\mathbf{X})}\) characterized by \(\nu(x) = \delta_x\) and \(\pi(x) = \delta_x\) for all \(x \in \mathbf{X}\) . If \(\iota\) is the injection of \(\mathbf{N}^{(\mathbf{X})}\) into \(\mathbf{Z}^{(\mathbf{X})}\) , the two homomorphisms \(\iota \circ \nu\) and \(\pi \circ \mu\) of \(\mathrm{Mo}(\mathbf{X})\) into \(\mathbf{Z}^{(\mathbf{X})}\) coincide on \(\mathbf{X}\) , whence \(\iota \circ \nu = \pi \circ \mu\) . The situation may be summarized by the following commutative diagram:

\[
\begin{array}{c} \mathrm{M} (\mathrm{X}) \xrightarrow {\lambda} \mathrm{Mo} (\mathrm{X}) \xrightarrow {\nu} \mathbf {N} ^ {(\mathrm{X})} \\ \mu \Biggl \downarrow \\ \mathrm{F} (\mathrm{X}) \xrightarrow {\pi} \mathbf {Z} ^ {(\mathrm{X})}. \end{array}
\]

The homomorphisms \(\lambda\) , \(\mu\) , \(\nu\) and \(\pi\) will be called canonical.

Let \(w\) be in \(\mathbf{M}(\mathbf{X})\) ; it is immediately shown by induction on \(l(w)\) that the length of the word \(\lambda(w)\) is equal to that of \(w\) . Moreover

\[
\nu (x _ {1} \dots x _ {n}) = \sum_ {i = 1} ^ {n} \delta_ {x _ {i}}\tag{21}
\]

for \(x_1, \ldots, x_n\) in X, whence \(|\nu(x_1 \ldots x_n)| = n\) by (13) and (14). In other words,

\[
\left| \nu (u) \right| = l (u) \quad (u \in \operatorname{Mo} (\mathbf {X})).\tag{22}
\]

PROPOSITION 11. The canonical homomorphism \(\nu\) of \(\operatorname{Mo}(\mathbf{X})\) into \(\mathbf{N}^{(\mathbf{X})}\) is surjective. Let \(w = x_1 \ldots x_n\) and \(w' = x_1' \ldots x_m'\) be two elements of \(\operatorname{Mo}(\mathbf{X})\) ; in order that \(\nu(w) = \nu(w')\) , it is necessary and sufficient that \(m = n\) and that there exist a permutation \(\sigma \in \mathfrak{S}_n\) with \(x_i' = x_{\sigma(i)}\) for \(1 \leqslant i \leqslant n\) .

The image of \(\nu\) is a submonoid I of \(\mathbf{N}^{(\mathbf{x})}\) containing the elements \(\delta_x\) (for \(x \in \mathbf{X}\) ). Formula (12) (no. 7) shows that \(\mathbf{N}^{(\mathbf{x})}\) is generated by the family \((\delta_x)_{x \in \mathbf{X}}\) , where \(\mathbf{I} = \mathbf{N}^{(\mathbf{x})}\) . Therefore \(\nu\) is surjective.

If \(m = n\) and \(x_{i}^{\prime} = x_{\sigma (i)}\) for \(\mathrm{l}\leqslant i\leqslant n\) , then

\[
\nu (w ^ {\prime}) = \sum_ {i = 1} ^ {n} \delta_ {x _ {i} ^ {\prime}} = \sum_ {i = 1} ^ {n} \delta_ {x _ {\sigma (i)}} = \sum_ {i = 1} ^ {n} \delta_ {x _ {i}} = \nu (w)
\]

by formula (21) and the commutativity theorem (§ 1, no. 5, Theorem 2).

Conversely, suppose that \(\nu(w)\) and \(\nu(w')\) are equal to the same element \(\alpha\) of \(\mathbf{N}^{(\mathbf{x})}\) ; by formula (22), \(n = |\alpha| = m\) . For all \(x \in \mathbf{X}\) , let \(\mathrm{I}_x\) (resp. \(\mathrm{I}'_x\) ) be the set of integers \(i\) such that \(1 \leqslant i \leqslant n\) and \(x_i = x\) (resp. \(x_i' = x\) ). Hence \((\mathrm{I}_x)_{x \in \mathbf{X}}\) and \((\mathrm{I}'_x)_{x \in \mathbf{X}}\) are partitions of the interval \([1, n]\) of \(\mathbf{N}\) ; further, the formula \(\alpha = \sum_{i=1}^{n} \delta_{x_i}\) shows that \(\alpha(x)\) is the cardinal of \(I_x\) ; similarly the formula \(\alpha = \sum_{i=1}^{\infty} \delta_{x'_i}\) shows that \(\alpha(x)\) is the cardinal of \(\mathbf{I}_x'\) . There thus exists a permutation \(\sigma\) of \([1, n]\) such that \(\sigma(\mathbf{I}_x') = \mathbf{I}_x\) for all \(x \in \mathbf{X}\) , that is \(x'_i = x_{\sigma(i)}\) for \(i = 1, \ldots, n\) .

Remark. Let S be a subset of X. Recall that we have identified M(S) with a submagma of M(X), Mo(S) with a submonoid of Mo(X) and \(\mathbf{N}^{(S)}\) with a submonoid of \(\mathbf{N}^{(X)}\) . Then

\[
\mathbf {M} (\mathrm{S}) = \lambda^ {- 1} (\mathbf {M o} (\mathrm{S})).\tag{23}
\]

Clearly \(\lambda (\mathbf{M}(\mathbf{S}))\subset \mathbf{Mo}(\mathbf{S})\) . Let \(w\in \lambda^{-1}(\mathbf{Mo}(\mathbf{S}))\) ; we show by induction on \(l(w)\) that \(w\in \mathbf{M}(\mathbf{S})\) . It is obvious if \(l(w) = 1\) . If \(l(w) > 1\) , we may write \(w = w_{1}w_{2}\) with \(w_{1},w_{2}\in \mathbf{M}(\mathbf{X}), l(w_{1}) <   l(w),l(w_{2}) <   l(w)\) . Then \(\lambda (w_1)\lambda (w_2)\in \mathbf{Mo}(\mathbf{S})\) hence \(\lambda (w_1)\in \mathbf{Mo}(\mathbf{S})\) and \(\lambda (w_2)\in \mathbf{Mo}(\mathbf{S})\) , whence \(w_{1}\in \mathbf{M}(\mathbf{S})\) and \(w_{2}\in \mathbf{M}(\mathbf{S})\) by the induction hypothesis and finally \(w\in \mathbf{M}(\mathbf{S})\)

Also

\[
\mathbf {M o} (\mathrm{S}) = \nu^ {- 1} (\mathbf {N} ^ {(\mathrm{S})}).\tag{24}
\]

This follows immediately from formula (21).

Further, \(\mathbf{N}^{(\mathbf{S})}\) is the set of elements of \(\mathbf{N}^{(\mathbf{X})}\) whose support is contained in S; if \((\mathrm{S}_i)_{i\in I}\) is a family of subsets of X of intersection S, then \(\mathbf{N}^{(\mathbf{S})} = \bigcap_{i\in I}\mathbf{N}^{(\mathbf{S}_i)}\) and formulae (23) and (24) imply

\[
\mathbf {M} (\mathrm{S}) = \bigcap_ {i \in \mathrm{I}} \mathbf {M} (\mathrm{S} _ {i}), \quad \mathbf {M o} (\mathrm{S}) = \bigcap_ {i \in \mathrm{I}} \mathbf {M o} (\mathrm{S} _ {i}).\tag{25}
\]

